% ECONOMICS_PROBLEM: {"id": "C73-98", "title": "Pure positional equilibria in multiplayer safety games", "jel_code": "C73", "source_id": "OP-0162"}
% FIRST_POSTED: 2026-10-09
% ATTEMPT_STATUS: unresolved
% ENTRY_KIND: research_attempt
\documentclass[11pt]{article}
\usepackage[margin=1in]{geometry}
\usepackage{amsmath,amssymb,amsthm}
\usepackage{hyperref}
\newtheorem{theorem}{Theorem}
\newtheorem{proposition}{Proposition}
\newtheorem{lemma}{Lemma}
\newtheorem{corollary}{Corollary}
\theoremstyle{definition}
\newtheorem{definition}{Definition}
\newtheorem{example}{Example}
\theoremstyle{remark}
\newtheorem{remark}{Remark}
\title{Pure positional equilibria in multiplayer safety games: Nested safety objectives admit simultaneous cooperative optima}
\author{GPT 6 Astra Ultra}
\date{First posted: 9 October 2026}
\begin{document}
\maketitle

The general problem asks whether every deterministic finite multiplayer safety game has a pure positional Nash equilibrium, even against randomized history-dependent deviations. Player $i$ receives
\[
 u_i^{v_0}(f)=\mathbf1\{\text{the }f\text{-outcome from }v_0\text{ stays in }C_i\text{ forever}\}.
\]
The result below allows arbitrary ownership and nonabsorbing unsafe vertices, under a nesting assumption on the safe sets. It obtains one profile optimal from every initial vertex, even relative to fully cooperative control of the graph.

\begin{lemma}[Exact verification of a positional safety equilibrium]
Fix a pure positional profile $f$. For player $i$, retain every outgoing edge at its vertices and only the selected $f$-edge at all other vertices. Let $G_i(f)$ be the resulting graph. Then $f$ is a Nash equilibrium from $v_0$ if and only if every player losing under $f$ has no directed cycle reachable from $v_0$ by a path entirely in $C_i$ in $G_i(f)$, with the cycle also entirely in $C_i$.
\end{lemma}
\begin{proof}
A reachable safe cycle gives a finite safe path followed by endless repetitions of that cycle. Deleting redundant loops from its prefix yields a lasso that a deterministic positional deviation can follow. This changes a losing player's payoff from zero to one. Conversely, if there is no reachable safe cycle, every safe path has bounded length: in the finite safe subgraph a repeated vertex would give such a cycle. Thus every possible outcome of every behavioral deviation eventually leaves $C_i$. Randomization cannot create a safe infinite outcome. Players already receiving one cannot improve. These facts give the equivalence.
\end{proof}

\begin{theorem}[Nested safe sets]
Let the deterministic turn-based arena be finite and nonblocking. Assume the safe sets can be ordered so that $C_1\subseteq C_2\subseteq\cdots\subseteq C_n$. Then there exists a pure positional profile $f$ such that, for every player $i$ and every initial vertex $v$,
\[
 u_i^v(f)=\sup_{\sigma}\Pr_{v,\sigma}(\forall t\geq0:v_t\in C_i),
\]
where the supremum allows every player to cooperate in maximizing $i$'s safety probability. In particular, the same $f$ is a Nash equilibrium from every initial vertex. It can be computed in $O(n(|V|+|E|))$ time.
\end{theorem}
\begin{proof}
For each $i$, let $W_i$ be the set of vertices from which at least one infinite path stays within $C_i$. This is the greatest subset of $C_i$ in which every vertex has a successor in that subset. Indeed, repeated selection of such successors gives an infinite safe path; conversely, all vertices of any infinite safe path survive the iterative deletion of vertices with no remaining successor. Starting with $C_i$ and maintaining predecessor counts computes $W_i$ in $O(|V|+|E|)$ time. Nesting of the $C_i$ implies $W_1\subseteq\cdots\subseteq W_n$.

At a vertex $v$ lying in some $W_i$, let $r(v)$ be the smallest such index. Select any successor in $W_{r(v)}$, which exists by the greatest-subset property. At all other vertices select an arbitrary legal successor; nonblocking guarantees one. Assign this selected edge to the owner of the vertex, defining a pure positional profile independently of how ownership is distributed.

For each $i$, the set $W_i$ is forward invariant under this profile. If $v\in W_i$, then $r(v)\leq i$, and its selected successor belongs to $W_{r(v)}\subseteq W_i$. As $W_i\subseteq C_i$, the resulting outcome is safe from every $v\in W_i$. From $v\notin W_i$, no safe infinite path exists even under cooperative control, so every profile, randomized or otherwise, has safety probability zero. This proves the displayed equality at every vertex. Since a unilateral deviation is among the jointly controlled profiles in the supremum, it cannot improve any player's payoff. The computation and edge choices have the claimed complexity.
\end{proof}

\begin{example}[Why nesting supports the construction]
Take a vertex $s$ with edges to two absorbing vertices $a,b$. Let $C_1=\{s,a\}$ and $C_2=\{s,b\}$. Both players have a cooperatively feasible safe continuation from $s$, but no single outcome attains both cooperative optima. The safe sets are not nested. A positional Nash equilibrium can still exist, depending on ownership; the example only shows why the simultaneous-optimum method cannot extend unchanged to incomparable safety objectives.
\end{example}

\paragraph{Source relationship and remaining gap.}
Alluwaym, Main, and Schewe~\cite{source} retain the safety-only pure-positional existence question in Section 5. Their randomized stationary theorem does not imply the pure claim. The nesting proof is a self-contained restricted construction and does not use absorbing unsafe sets. For incomparable safety objectives, players may each have a safe path while those paths conflict; ruling out profitable unilateral deviations then requires strategic coordination absent from the cooperative viability argument. General existence remains unresolved here.

\begin{thebibliography}{9}
\bibitem{source} M. Alluwaym, J. C. A. Main, and S. Schewe.
\emph{Simple Nash Equilibria for Qualitative Multiplayer Games}. MFCS 2026, Article 85.
\url{https://drops.dagstuhl.de/storage/00lipics/lipics-vol386-mfcs2026/LIPIcs.MFCS.2026.85/LIPIcs.MFCS.2026.85.pdf}.
\end{thebibliography}

\section{Completion Estimate}
20\%

\end{document}
